% !TEX root = ../main.tex
\chapter{Literature Review and Theoretical Framework}
\label{ch:literature}

The research problem set out in Chapter~\ref{ch:introduction} is that transfer-graph reputation scores give no role to spending authorization, an explicit act of trust that every ERC-20 token records, and that their graphs grow with every payment. The objectives O1--O5 and research questions RQ1--RQ5 follow from it. To compare EndorseRank with Adaptive Weighted PageRank (AWP) on this problem, we run each method's graph through PageRank and check the resulting rankings with rank correlations. The framework combines hyperlink--citation propagation, domain alignment and computational complexity. The gap that remains is taken up in Chapters~\ref{ch:methodology} and~\ref{ch:results}.

\dissertationSubheading[sec:blockchain-defi-foundations]{Blockchain and DeFi foundations}

A public blockchain records transactions in hashed blocks, and each block is secured by the ones added after it \cend{5}. Once a majority of the participating nodes has accepted a block, rewriting it is costly. That cost allows parties to settle on the ledger without each of them needing a known legal identity (\emph{trust-minimized} settlement).

Bitcoin built a peer-to-peer electronic cash network on this kind of ledger \cend{5}. Ethereum keeps the append-only ledger but lets an account hold a program alongside its balance. A transaction can then call one of the program's functions and change the contract's state \cend{7}, and these functions run on the Ethereum Virtual Machine (EVM) \cend{36}. Programs stored on the ledger in this way are called smart contracts; any honest node that replays a function call gets the same output \cend{6}. As a result, collateralization ratios, liquidation conditions and fee schedules can be enforced by the contract instead of by a trusted third party \cend{1}.

Collectively, these contracts make up Decentralized Finance (DeFi). Between them they do the work of exchanges, lending platforms, derivatives markets and asset managers; Werner et al.\cend{8} define the protocol classes, and Sch\"{a}r\cend{1} describes how protocols are layered on one another. Because one protocol can interact with another, new financial instruments can be composed from existing ones \cend{37}. The same property lets a failure spread across several protocols. Successful attacks have included flash-loan exploits \cend{38} and price-oracle manipulation \cend{39}. Transaction histories on these chains are public and machine-readable, yet they record actions, not legal identities \cend{2}.

\dissertationSubsubheading[sub:pseudonymity-credit]{Pseudonymity and the credit problem}

Conventional credit scoring ties a credit history to a legal entity. On a public blockchain the basic unit of analysis is instead an address (a wallet), which may stand for a legal entity, a trust, an algorithmic agent or a smart contract. One entity can own several addresses \cend{11}, and statistical clustering can recover ownership only in part \cend{40}. On-chain lending therefore requires borrowers to post collateral worth more than the amount they borrow; if the collateral value falls below a set threshold, a liquidator may repay part of the debt and receive the corresponding collateral at a discount \cend{41}. The arrangement is secure, but it ties up capital. An on-chain credit score tries to quantify from on-chain behavior alone how likely an actor is to repay. Recent proposals draw on prospect theory \cend{13}, on machine learning applied to lending data \cend{14}, and on models of credit risk spreading through account networks \cend{12}.

\dissertationSubheading[sec:ethereum-l2]{Ethereum, Layer-2 rollups, and Arbitrum One}

Most DeFi applications have been deployed on the Ethereum mainnet, historically the main venue for DeFi innovation and liquidity, but growing adoption there has raised gas costs and constrained throughput \cend{42}. Layer-2 (L2) protocols try to keep Ethereum's security guarantees while making transactions cheaper. Some execute transactions off-chain, while others batch them and post a succinct summary of the execution to the Ethereum mainnet \cend{43}. Thibault et al.\cend{44} organize L2 scaling approaches into a taxonomy. Optimistic L2s such as Arbitrum treat each submitted batch of transactions as valid unless evidence against it is presented within a dispute period \cend{45}.

Arbitrum One is an optimistic L2 of this kind and can execute EVM-based smart contracts. It has become a popular platform for DeFi applications, among them GMX, an exchange for perpetual swaps \cend{46}. We use Arbitrum One because the public BigQuery dataset includes both its ERC-20 \texttt{Approval} and \texttt{Transfer} events, and because its fees are low enough that a six-month period contains dense activity. We make no claim that our results generalize to other L2s or to Ethereum mainnet without further validation. We do believe, however, that the proposed framework (approval and transfer graphs, construct checks, and rank correlation) applies to any chain that emits the same logs.

\dissertationSubheading[sec:erc20-allowance-lifecycle]{ERC-20 tokens and the allowance lifecycle}

The ERC-20 standard defines a common interface for fungible tokens, made up of \texttt{transfer}, \texttt{approve}, \texttt{transferFrom}, \texttt{allowance} and the associated events \cend{18}. After the owner calls \texttt{approve(spender, amount)}, the token contract records that the spender may move up to \texttt{amount} out of the owner's balance. The spender can then call \texttt{transferFrom} to carry out a token transfer. With EIP-2612 \texttt{permit}, an allowance can be updated through an off-chain signature that settles in a single transaction \cend{19}. The user experience improves, and the underlying semantics of authorization are unchanged.

Figure~\ref{fig:allowance-sequence} traces an ERC-20 allowance through its lifecycle. An allowance is an authorization and involves no actual transfer. Token owners can grant allowances to routers, vaults or trading contracts that they trust to act within set limits. Revoking (approving zero) and updating (setting a new non-zero allowance) both change the allowance currently in force. EndorseRank accounts for this by taking the most recent allowance value for each owner--spender pair as the edge weight, so a newer allowance replaces the previous one outright; old values do not decay gradually over time.

\begin{figure}[htbp]
\centering
\begin{tikzpicture}[
  font=\small,
  actor/.style={draw, rounded corners, minimum width=2.4cm, minimum height=0.75cm, align=center},
  msg/.style={-{Stealth[length=2.5mm]}, thick},
  note/.style={font=\scriptsize\itshape, text=black!70, align=left}
]
% Actors
\node[actor] (owner) at (0,0) {Owner};
\node[actor] (token) at (5.5,0) {ERC-20 token};
\node[actor] (spender) at (11,0) {Spender};

% Lifelines (vertical dashed lines)
\draw[dashed, gray] (owner.south) -- ++(0,-6.6);
\draw[dashed, gray] (token.south) -- ++(0,-6.6);
\draw[dashed, gray] (spender.south) -- ++(0,-6.6);

% Message 1: approve (owner -> token), horizontal span 5.5cm
\draw[msg] (0,-1.2) -- node[above,font=\scriptsize] {\texttt{approve(spender,$a$)}} (5.5,-1.2);
\node[note] at (2.75,-1.7) {emit \texttt{Approval}; store allowance};

% Alternative to message 1: EIP-2612 permit, submitted by the spender (or a relayer) with the owner's signature
\draw[msg, densely dotted] (11,-2.6) -- node[above,font=\scriptsize] {\texttt{permit(owner,spender,$a$,sig)}} (5.5,-2.6);
\node[note] at (8.25,-3.1) {alternative to \texttt{approve}: owner signs off-chain;\\ emit \texttt{Approval}; store allowance};

% Message 2: transferFrom (spender -> token), horizontal span 5.5cm
\draw[msg] (11,-4.2) -- node[above,font=\scriptsize] {\texttt{transferFrom(owner,to,$v\leq a$)}} (5.5,-4.2);
\node[note] at (8.25,-4.7) {move tokens; reduce allowance};

% Message 3: revoke (owner -> token), dashed
\draw[msg, dashed] (0,-5.8) -- node[above,font=\scriptsize] {\texttt{approve(spender,$0$)} (revoke)} (5.5,-5.8);
\end{tikzpicture}
\caption[{ERC-20 allowance lifecycle: approve (or permit), optional \texttt{transferFrom}, and revoke. An EndorseRank edge sums, over tokens, the latest in-window approval of its owner--spender pair, weighted by its age and amount.}]{ERC-20 allowance lifecycle: \texttt{approve}, or its signature-based alternative \texttt{permit} (dotted), an optional \texttt{transferFrom}, and revoke. An EndorseRank edge sums, over tokens, one term for the most recent \texttt{Approval} event of its owner--spender pair in the observation window, and the term weights that approval by its age and amount as AWP weights a transfer. A later approval replaces an earlier one.}
\label{fig:allowance-sequence}
\end{figure}

In the available data, allowances are sparse compared with transfers. A single approval can cover many transfers, and many wallets have never granted an approval to any external account. The sparsity is a computational advantage, and it also marks a semantic difference, since EndorseRank quantifies authorization to spend whereas a transfer graph records payment flows.

\dissertationSubheading[sec:graph-reputation]{Graph-based on-chain reputation}

PageRank and its variants rank nodes by the influence that flows to them along directed edges \cend{15}. Two recent studies apply this approach to blockchain transaction graphs, and the two should be distinguished. Do and Do\cend{4} run PageRank and HodgeRank on Ethereum transaction graphs to measure social credit; their main contribution is a definition of connectivity between accounts for credit scoring. Do, Do and Nguyen\cend{3} propose Adaptive Weighted PageRank (AWP), the transfer-graph method that this thesis evaluates beside EndorseRank (Section~\ref{sec:awp-in-depth}). Lin et al.\cend{12} argue that transaction connectivity can miss risk-relevant structure unless risk is propagated through the network, and Wu et al.\cend{24} reach a similar conclusion for phishing detection, where network embeddings of transaction graphs outperform per-account features. Surveys of blockchain-based reputation systems \cend{25} and of smart-contract reputation designs \cend{26} find that most deployments still take explicit ratings or transfer volume as their trust signal, and a survey of knowledge discovery in cryptocurrency transaction data \cend{47} names degree- and value-based centralities as the main descriptors of wallets. Packin and Lev-Aretz\cend{2} stress that such systems must be auditable even though their users are pseudonymous.

Outside the research literature, trust in deployed DeFi and web3 systems is usually established another way. An address builds standing through activity that anyone can inspect on the ledger, and tooling bound to the address summarizes it. The most widely used tool, Human Passport (formerly Gitcoin Passport), collects verifiable credentials, called stamps, for checks such as government identity, biometrics, vouching and earlier web2 and web3 activity. Each stamp carries a weight that reflects how hard it is to forge, and an address passes when the weighted sum reaches a threshold \cend{64}. Passport also classifies EVM addresses on several networks, Arbitrum among them, as human or Sybil from their on-chain activity \cend{64}. Its purpose is Sybil resistance, so it combines on-chain with off-chain evidence and works as a gate and not as a graded ordering. A review of proof-of-personhood protocols finds that they trade Sybil resistance against self-sovereignty and privacy \cend{65}, and a recent security review names credential farming, the bulk acquisition of cheap stamps, as the main residual weakness of Passport \cend{66}. Such tools answer whether an address is controlled by a distinct person. EndorseRank answers a different question, which addresses other addresses authorize to spend their tokens, and it orders them from on-chain records alone (Table~\ref{tab:method-comparison-lit}).

Graph-based methods are attractive because they (i) use only public information, (ii) produce continuous scores that suit ranking, and (iii) admit complexity analysis by number of edges. They are, however, vulnerable to Sybil attacks \cend{11} and wash trading, and they suffer from the mismatch between graph centrality and actual default risk. PageRank scores are not credit scores in the regulatory sense, and the aim here is only to compare allowance-based and transfer-based rankings rigorously against the same set of proxies.

\dissertationSubheading[sec:pagerank-formalism]{PageRank formalism}

Consider a directed graph $G=(V,E)$ with $|V|$ nodes, and write $w_{ij}>0$ for the weight of the edge from $i$ to $j$. That weight stands for influence from $i$ to $j$, in either the endorsement or the transfer sense introduced above. Let $W$ be the corresponding weighted adjacency matrix, and let $o_i=\sum_j w_{ij}$ be the out-strength of node $i$. The row-stochastic transition matrix $P$ is then
\begin{equation}
P_{ij} =
\begin{cases}
w_{ij}/o_i & \text{if } o_i>0,\\
1/|V| & \text{if } o_i=0
\end{cases}
\label{eq:transition}
\end{equation}
(dangling nodes spread their mass evenly). Given a damping parameter $d\in(0,1)$, the Google matrix is
\begin{equation}
G_{\mathrm{PR}} = d\, P^{\top} + \frac{1-d}{|V|}\,\mathbf{1}\mathbf{1}^{\top},
\label{eq:google-matrix}
\end{equation}
and the PageRank vector $\mathbf{r}$ satisfies $\mathbf{r}=G_{\mathrm{PR}}\mathbf{r}$ with $\mathbf{r}\geq 0$ and $\sum_i r_i=1$ \cend{15,23}. Element by element, with the dangling rows of $P$ written as a separate middle term, this becomes
\begin{equation}
r_j = d\sum_{i:(i\to j)\in E} r_i \frac{w_{ij}}{o_i} + \frac{d}{|V|}\sum_{i:\,o_i=0} r_i + \frac{1-d}{|V|}.
\label{eq:pagerank-node}
\end{equation}

\dissertationSubsubheading[sub:random-surfer]{Random-surfer interpretation}

Equation~\eqref{eq:pagerank-node} can be read probabilistically. At each step a surfer at a node with out-edges follows one of them with probability $d$, choosing it in proportion to its weight, or teleports to a uniformly chosen node with probability $1-d$ (Figure~\ref{fig:random-surfer}). At a node without out-edges the surfer always jumps to a uniformly chosen node; this is the dangling rule of Equation~\eqref{eq:transition}. The stationary distribution then ranks nodes by how often they are visited in the long run. Damping keeps rank sinks (groups of nodes with no edge leaving them) from absorbing all the mass. With $d<1$ and uniform teleportation, every entry of $G_{\mathrm{PR}}$ is positive, so $G_{\mathrm{PR}}$ is primitive and its Perron vector is unique \cend{23}. Appendix~\ref{sec:pr-existence} gives the argument.

\begin{figure}[htbp]
\centering
\begin{tikzpicture}[
  font=\small,
  vertex/.style={draw, circle, minimum size=0.95cm},
  arrow/.style={-{Stealth[length=2.5mm]}, thick},
  teleport/.style={arrow, dashed}
]
\node[vertex] (a) at (0,0) {$A$};
\node[vertex] (b) at (3.2,1.5) {$B$};
\node[vertex] (c) at (3.2,-1.5) {$C$};
\node[vertex] (d) at (6.4,0) {$D$};
% Graph transitions: A->B, A->C, B->D, C->D (same topology as the worked example; D is dangling)
\draw[arrow] (a) -- node[above left, font=\footnotesize] {$2$} (b);
\draw[arrow] (a) -- node[below left, font=\footnotesize] {$1$} (c);
\draw[arrow] (b) -- node[below left=-1pt, font=\footnotesize] {$3$} (d);
\draw[arrow] (c) -- node[above left=-1pt, font=\footnotesize] {$1$} (d);
% Uniform teleportation, drawn from D only: one dashed arrow to every node, including D itself.
% D->A runs through the gap between B and C, and D->B / D->C bend outward, so no dashed arrow crosses a solid edge.
\draw[teleport] (d) -- (a);
\draw[teleport] (d) to[bend right=32] (b);
\draw[teleport] (d) to[bend left=32] (c);
\draw[teleport] (d) to[loop right, min distance=12mm] (d);
\end{tikzpicture}
\caption[{Random-surfer intuition for PageRank on the four-node graph of Section~\ref{sec:pagerank-worked-example}. Solid arrows are graph transitions, followed with total probability $d$. Dashed arrows are uniform teleportation, probability $(1-d)/n$ to every node including the current one; the dashed fan is drawn from $D$ only but leaves every node.}]{Random-surfer intuition for PageRank on the four-node graph of Section~\ref{sec:pagerank-worked-example}; edge labels are the allowance weights. Solid arrows are graph transitions: with total probability $d$ the surfer follows one of them, chosen in proportion to its weight. Dashed arrows show uniform teleportation, which reaches every node (the current one included) with probability $(1-d)/|V|$; the dashed fan is drawn only from $D$, although every node has one. $D$ has no outgoing edge, so under the dangling rule of Equation~\eqref{eq:transition} a surfer at $D$ always jumps to a uniformly chosen node, each with probability $1/|V|=0.25$. That rule, and not damping, keeps $D$ from absorbing the chain.}
\label{fig:random-surfer}
\end{figure}

\dissertationSubsubheading[sub:power-iteration]{Power iteration and complexity}

In practice, $\mathbf{r}$ is usually obtained by power iteration from the uniform vector $\mathbf{r}^{(0)}=\mathbf{1}/|V|$,
\begin{equation}
\mathbf{r}^{(t+1)} = G_{\mathrm{PR}}\mathbf{r}^{(t)},
\label{eq:power-iteration}
\end{equation}
until $\|\mathbf{r}^{(t+1)}-\mathbf{r}^{(t)}\|_1 < \varepsilon$ or until the iteration cap $I_{\max}$ is reached. On a sparse graph, each iteration costs $O(|E|)$ arithmetic operations \cend{48}. A method that generates fewer edges, for instance by using latest-allowance snapshots instead of full transfer histories, therefore lowers the cost per iteration and sometimes the memory required. Hypothesis H3 is motivated by this complexity fact.

By default this work sets damping to the conventional $d=0.85$ \cend{15} and uses two study-specific solver settings, a convergence tolerance $\varepsilon=10^{-8}$ and a maximum of $I_{\max}=300$ iterations.\footnote{The tolerance and the iteration cap were chosen to keep run-times reproducible; they are not put forward as part of the classical PageRank formulation.} Damping is varied over $d\in\{0.75,0.85,0.95\}$ in the robustness section.

\dissertationSubheading[sec:reputation-systems]{Reputation systems lineage}

PageRank belongs to a wider family of spectral and propagation-based reputation mechanisms. Kleinberg's HITS algorithm assigns each page two scores, one as a hub and one as an authority \cend{49}. EigenTrust extends the idea to peer-to-peer networks and normalizes local trust, so a malicious node cannot raise its own score at will \cend{16}. To suppress web spam, TrustRank starts from a seed set of trusted nodes and propagates their trust outward \cend{17}. Reputation is relational in all of these methods. Where trust flows only from a vetted seed set, as in TrustRank, a newcomer cannot raise its score without endorsements from nodes that already hold reputation. With uniform teleportation, which is the choice made here so that EndorseRank matches AWP, that protection is weaker, and Section~\ref{sec:sybil-cost-model} works out by how much.

On-chain, the Sybil attack is still the main concern \cend{11}. New wallets cost little to create, whereas wallets that receive allowances or transfers from distinct counterparties cost more. The Sybil-adjusted stability proxies used in this work (inbound counterparty ratio, tenure, active months) partly capture that diversity of interaction, but they do not prevent a malicious actor from manipulating them. Chapter~\ref{ch:discussion} presents a cost model for Sybils in allowance graphs.

Machine-learning credit models that use DeFi features \cend{14} are trained to predict labeled defaults, and they give up interpretability to do so. EndorseRank and AWP keep to interpretable, audit-friendly propagation rules \cend{2}.

\dissertationSubheading[sec:allowance-endorsement]{ERC-20 allowances as endorsement}

Through the ERC-20 \texttt{approve} function and EIP-2612 \texttt{permit}, users state explicitly how much a spender may move on their behalf \cend{18,19}. EndorseRank takes the latest allowance from an owner to a spender on each token as one endorsement, weights it as AWP weights a transfer, and runs PageRank on the resulting directed graph. Newer approvals supersede older ones, so each pair and token contributes a single event.

Concretely, for each owner--spender pair the latest in-window approval on each token contributes its time weight times its bounded value weight, as AWP weights a transfer; a pair whose latest value is zero (a revocation) is dropped, and the terms are summed over tokens (Equation~\eqref{eq:approve-weight} in Section~\ref{sub:endorserank-graph}). The resulting graph $G_{\text{approve}}$ is usually much sparser than the transfer graph over the same wallets and window, and PageRank on it yields the EndorseRank scores $\mathbf{s}_{\mathrm{ER}}$.

Treating allowances as endorsements is a deliberate choice, and an imprecise one. Users approve routers because it is convenient, not necessarily to make a social statement about whether the router is trustworthy. Contract addresses and EOAs coexist as well. Even so, an allowance signals approval more explicitly than a received transfer, which could come from market making, wash trading \cend{29}, airdrop farming \cend{32} or routing without any intent to endorse.

\dissertationSubheading[sub:gf-pr]{What this research does not claim}

The research proposal planned to validate both methods against GMX~V2 trading success and to bound that validation with an outcome-based reference, GF-PR: PageRank on a star graph built from realized profit and loss. GF-PR is circular. It is computed from the same profit-and-loss stream that defines the trading-success proxies, so its near-perfect agreement with the realized-gain proxy is the correlation of a variable with a monotone transformation of itself, and the research question and claim that depended on it were withdrawn with it. The same-window trading proxies remain only as archived families (Table~\ref{tab:alignment-summary}). Trading outcomes return out of window, in an exploratory run on the matched traders (Section~\ref{sub:trader-results}) and as the liquidation-free close rate of the registered June--August window (Section~\ref{sub:fresh-holdout}). No score in this thesis is offered as a predictor of trading profit, loss avoidance or liquidation.

\dissertationSubheading[sec:rank-correlation-theory]{Rank correlation theory}

Reputation scores and validation proxies are both continuous variables, and neither has a natural boundary for classification. Following the AWP validation design\cend{3} and classical psychometrics \cend{22}, this research measures two kinds of rank agreement, monotonic and pairwise. Bootstrap confidence intervals quantify the sampling uncertainty of every reported coefficient \cend{34}.

Spearman's rank correlation \cend{20} is the Pearson correlation of the ranks. For paired ranks $(R_i,S_i)$ of $n$ items without ties, it reduces to
\begin{equation}
\rho = 1 - \frac{6\sum_{i=1}^n (R_i-S_i)^2}{n(n^2-1)}.
\label{eq:spearman}
\end{equation}
This no-ties form does not hold when ties occur. Spearman's $\rho$ is then the Pearson correlation of average ranks (mid-ranks), in which tied items share the mean of the positions they occupy, and every reported $\rho$ is computed this way. In both forms $\rho$ measures monotonic association.

Kendall's $\tau$ \cend{21} rests on counts of concordant and discordant pairs. Of the $n_0=n(n-1)/2$ pairs $(i,j)$ with $i<j$, $n_c$ are concordant (both rankings order them the same way) and $n_d$ are discordant (the rankings order them oppositely); a pair tied in either ranking is neither. Let $n_1=\sum_k t_k(t_k-1)/2$, where $t_k$ is the size of the $k$-th group of tied values in the first ranking, and define $n_2$ in the same way for the second ranking. The $\tau$-b variant is then
\begin{equation}
\tau_b = \frac{n_c - n_d}{\sqrt{(n_0-n_1)(n_0-n_2)}},
\label{eq:kendall}
\end{equation}
and without the tie corrections it becomes $\tau_a=(n_c-n_d)/n_0$. Every $\tau$ reported in this thesis is $\tau_b$, computed on the raw scores and proxies. The difference matters for heavily tied scores, and EndorseRank on the matched cohort is one (Section~\ref{sec:rank-divergence}). For agreement about pairwise orderings, Kendall's $\tau$ is usually the easier coefficient to interpret, and Chapter~\ref{ch:results} uses it as the main family-level summary. Its values are compared with those of other methods on the same cohort; no uncited absolute benchmark is used.

\dissertationSubheading[sec:theoretical-framework]{Theoretical framework}

This study compares EndorseRank and AWP on reputation structure (H1), construct alignment (H2) and computational efficiency (H3), and then on prediction out of window (RQ4) and on coupling the two layers (RQ5). Three bodies of theory frame the first part: hyperlink--citation propagation, construct (domain) alignment, and computational complexity. Two further arguments bear on the second part.

\begin{itemize}
\item Hyperlink--citation semantics: In the original PageRank algorithm \cend{15}, an inbound hyperlink counts as a citation. Do and Do\cend{4} apply this reasoning to Ethereum transfers, and AWP\cend{3} weights each transfer edge by the age and the value of its transfers, through a logistic decay in age and a bounded transform of value (Section~\ref{sub:awp-graph}). EndorseRank extends citation-based propagation to ERC-20 approve and permit edges. An allowance is an explicit, on-chain delegation of spending authority that remains in force until it is revoked or replaced \cend{18,19}, whereas a transfer may reflect trading or routing without any intent to trust. Graph-based risk models \cend{12} show that network structure can carry information beyond raw activity counts, which argues for a directed endorsement graph alongside transfer volume. Whether propagation adds anything over the count on a particular label is an empirical matter, and the holdouts test it directly.

\item Domain alignment framework: Measurement theory uses construct alignment to ask how closely an indicator follows the construct it is meant to track \cend{22}. Messick\cend{33} widens the idea into one unified judgment about what a score means and what using it implies. With rank-based reputation scores, Spearman's $\rho$ and Kendall's $\tau$ measure monotonic and pairwise agreement in ordering and need no classification threshold, as in the AWP validation protocol\cend{3}. Graph coherence is checked with contemporaneous transfer and allowance statistics. The out-of-window checks are temporal holdouts: new approvals and new transfer senders for spenders, and, in the registered June--August window, the liquidation-free close rate for traders.

\item Computational complexity: Each PageRank iteration costs $O(|E|)$ \cend{15}. Allowance graphs are sparser than transfer graphs, with fewer edges $|E|$, so an iteration needs less time and memory, and in practice convergence often takes fewer iterations too. This argument underlies H3. Both methods use the same solver, so any difference in runtime is an expected effect of edge sparsity and does not count as an algorithmic contribution.

\item Local counts as a baseline: A wallet that already has many counterparties tends to gain more of them. Where that holds, the raw degree is a strong predictor of future arrivals, and a propagated score can beat it only when who the endorsers are matters more than how many there are. For this reason RQ4 sets each PageRank against its own raw degree at the freeze date.

\item Layer coupling: Multilayer PageRank treats several relations as layers over one node set \cend{62,63}. Mixing the layers node by node, as C-PR does, leaves a question open for layers of very different density: a node with no edge in the sparse layer follows the dense one, so the signal of the sparse layer may or may not survive. RQ5 therefore asks two things: whether the mixture improves on the transfer walk, and whether it beats both layers.
\end{itemize}

The conceptual model in Figure~\ref{fig:conceptual-model} links the scores of the thesis to H1--H3 and RQ1--RQ4. RQ5 is the dashed box on the right. It takes input from H2/RQ2 and from RQ4 and asks whether the two constructs combine.

\begin{figure}[htbp]
\centering
% !TEX root = ../main.tex
\begin{tikzpicture}[
  font=\small,
  box/.style={draw, rounded corners, align=center, minimum width=10.5cm, minimum height=1.0cm},
  smallbox/.style={draw, rounded corners, align=center, minimum width=3.4cm, minimum height=1.05cm},
  refbox/.style={draw, dashed, rounded corners, align=center, minimum width=3.4cm, minimum height=0.95cm},
  arrow/.style={-{Stealth[length=2.5mm]}, thick}
]
\node[box] (center) at (0,3.2) {한 지갑 코호트에서 구한 허락 점수와 송금 점수};

\node[smallbox] (rep) at (-5,1.2) {평판\\구조};
\node[smallbox] (risk) at (0,1.2) {구성개념\\점검};
\node[smallbox] (eff) at (5,1.2) {계산\\비용};

\node[smallbox] (h1) at (-5,-1.0) {H1 / RQ1\\엔도스랭크--AWP\\갈라짐};
\node[smallbox] (h2) at (0,-1.0) {H2 / RQ2\\각 구성개념과\\맞는 정도};
\node[smallbox] (h3) at (5,-1.0) {H3 / RQ3\\점수마다\\드는 비용};

\node[refbox] (hold) at (0,-3.0) {RQ4: 시간 홀드아웃\\점수 대 원시 차수};
\node[refbox] (rq5) at (5,-3.0) {RQ5: 결합 연산자 C-PR\\화살표 두 종류, 라벨 두 개};

\draw[arrow] (center.south -| rep.north) -- (rep.north);
\draw[arrow] (center.south) -- (risk.north);
\draw[arrow] (center.south -| eff.north) -- (eff.north);

\draw[arrow] (rep.south) -- (h1.north);
\draw[arrow] (risk.south) -- (h2.north);
\draw[arrow] (eff.south) -- (h3.north);

\draw[arrow, dashed] (h2.south) -- (hold.north);
\draw[arrow, dashed] (h2.south east) -- (rq5.north west);
\draw[arrow, dashed] (hold.east) -- (rq5.west);
\end{tikzpicture}

\caption[{Conceptual model: the allowance and transfer scores mapped to reputation structure (H1/RQ1), construct checks (H2/RQ2), computational cost (H3/RQ3), and the temporal holdout of future new approvals and senders against the raw degrees (RQ4).}]{Conceptual model linking the allowance and transfer scores of one wallet cohort to reputation structure (H1/RQ1), construct checks (H2/RQ2), computational cost (H3/RQ3) and the temporal holdout, in which each PageRank is set against the raw degree it smooths (RQ4). The dashed box on the right is RQ5, the coupled operator C-PR.}
\label{fig:conceptual-model}
\end{figure}

Only EndorseRank and AWP enter H1. H2/RQ2 and H3/RQ3 concern every score, and RQ4 is a time-split test on the spender cohort with future new owner--spender approvals and new transfer senders as its labels and the raw degrees at the freeze as its baselines; it does not rank trading outcomes.

\dissertationSubheading[sub:terminology]{Terminology}

The formal constructs are defined in Chapter~\ref{ch:introduction}. The terms below are used when results are interpreted:

\begin{itemize}
\item Single-layer methods: EndorseRank (allowance graph) and AWP (transfer graph), the two PageRank scores that walk on one edge type. The coupled walks of Section~\ref{sub:coupled-graph} and the raw degrees are evaluated beside them.
\item Primary graph construct: the graph that a method is \emph{designed} to rank, for example $G_{\text{approve}}$ for EndorseRank, as distinct from the checks applied after scoring.
\item Contemporaneous check: a local transfer or allowance statistic taken from the same window as the score \cend{22}.
\item Temporal holdout: scores are frozen at $t_1$ (28~February~2026), and new owner--spender approval pairs and new transfer senders are counted up to $t_2$ (31~May~2026); a registered second window freezes scores on 31~May~2026 and counts labels from June to August~2026 (Section~\ref{sub:fresh-holdout}).
\item Matched wallet cohort ($n=5{,}521$), spender holdout cohort ($n=1{,}335$) and expanded wallet pool ($N=99{,}080$): the three address sets of Section~\ref{sub:cohorts}.
\end{itemize}

\dissertationSubheading[sec:domain-alignment-metrics]{Domain alignment metrics}

For rank-based scores, construct alignment is measured by Spearman's $\rho$ and Kendall's $\tau$ between scores and check rankings, following the AWP validation protocol\cend{3} and Cronbach and Meehl\cend{22}, on contemporaneous transfer and allowance checks and on the two labels of the temporal holdout; Chapter~\ref{ch:methodology} gives the operational definitions. The literature reviewed so far thus supports three design choices: PageRank as an interpretable propagation backbone, allowances as endorsement edges separate from transfers, and a check design based on rank correlation plus a time split.

\dissertationSubheading[sec:pagerank-worked-example]{Worked example: PageRank on a four-node endorsement graph}

Take four wallets $\{A,B,C,D\}$ joined by the latest-allowance edges
\[
A\to B\ (2),\quad A\to C\ (1),\quad B\to D\ (3),\quad C\to D\ (1),
\]
with damping $d=0.85$. The out-strengths are $o_A=3$, $o_B=3$, $o_C=1$ and $o_D=0$ (the last is dangling). The edges $B\to D$ and $C\to D$ are the only out-edges of their nodes, so each is followed with probability one and its weight has no effect on the scores. Teleportation adds $(1-d)/4=0.0375$ to every node at every step, the first one included. One power-iteration step from the uniform vector $\mathbf{r}^{(0)}=\mathbf{1}/4$ gives $\mathbf{r}^{(1)}=(0.0906,\,0.2323,\,0.1615,\,0.5156)$ for $(A,B,C,D)$. The mass already tilts toward $D$, since both $B$ and $C$ endorse $D$, while $A$ splits its mass between $B$ and $C$. The iteration converges to the stationary vector $\mathbf{r}=(0.1375,\,0.2154,\,0.1765,\,0.4706)$, so $D$ comes first, $B$ and $C$ follow, and $A$ is last. No wallet endorses $A$. Apart from teleportation, its whole score comes from the dangling rule, which spreads $D$'s mass evenly over the four nodes ($d\,r_D/4=0.1000$). The example shows that a node gains score by \emph{receiving} endorsements from nodes that are endorsed in turn, which is precisely the citation logic of PageRank \cend{15,23}.

Suppose the same wallets were linked by transfers instead of allowances. The transfers would in general join other pairs, or the same pairs with other per-pair totals. That gives another transition matrix $P$, and possibly a different ranking of exactly the same nodes. Hypothesis H1 is this situation in miniature: EndorseRank and AWP need not agree.

\dissertationSubheading[sec:defi-risk-taxonomy]{DeFi risk taxonomy and where reputation fits}

Werner et al.\cend{8} classify DeFi risk by protocol layer; Zhou et al.\cend{27} group it by the kinds of attack seen in practice. Bugs in contract code and misconfigured upgrade keys are smart-contract risks \cend{50}. A stale price, or a feed under someone's control, is an oracle risk \cend{39}. Market risks \cend{51} cover volatility and liquidation cascades, including stress in stablecoin and funding markets \cend{52}, and the market crash of March 2020 showed how these two layers can interact \cend{9}. Token-weighted voting, which can be modeled, falls under governance risk \cend{53}. We study a narrower layer, \emph{counterparty} or \emph{behavioral} risk: whether an address's past behavior makes it a safe counterparty, as a trader that handles leveraged exposure responsibly or as a contract that many owners trust to spend their tokens.

On-chain reputation scores cover only part of the behavioral risk layer and are no substitute for collateral, auditing, or oracle design. Their purpose is to compress a system's relational structure (who authorizes whom, who pays whom) into an ordinal measure that can be audited and reproduced with little effort \cend{2}. Within behavioral risk, EndorseRank looks at authorization and AWP at payment flows. Neither tries to predict exploits or governance attacks.

Finally, fake activity makes behavioral risk harder to measure. Wash trading inflates volume on decentralized exchanges \cend{28} and, on a larger scale, on unregulated centralized exchanges \cend{29}. Front-running and related forms of extractable value create transfers that follow ordering strategies instead of relationships \cend{54}; Qin et al.\cend{30} estimate how much of that value is captured on Ethereum. Scam tokens on Uniswap \cend{31} and cross-exchange arbitrage trades \cend{55} likewise add transactions without forming sustained relationships. The transfer graph is particularly sensitive to fake volume, while the allowance graph is vulnerable to approval farming. Time-based splits and Sybil detection solve the problem only in part.

\dissertationSubheading[sec:related-measurement]{Related measurement traditions}

Outside blockchains, reputation systems have been studied in e-commerce \cend{56}, in peer-to-peer networks \cend{16}, and in work on web spam \cend{17}. Four insights from this work apply here: (i) reputation is a property of relations; (ii) when identities are easy to create, Sybil attacks become possible \cend{11}; (iii) seeding and normalization limit the effect of inflation; and (iv) evaluation needs external benchmarks, because internal coherence alone is not enough \cend{22}. EndorseRank takes up (i) directly and treats (ii) as a threat. It does not take up (iii): to stay comparable with AWP it keeps uniform teleportation and no seed set, and Chapter~\ref{ch:discussion} measures what that costs. On (iv) it goes part of the way. The contemporaneous checks, measured by Spearman's $\rho$ \cend{20} and Kendall's $\tau$ \cend{21}, test coherence, while the temporal holdouts test prediction out of window, on the same constructs for spenders and on liquidation for traders.

Personalized and topic-sensitive PageRank \cend{57,58} bias the teleportation vector toward a chosen group of nodes, the mechanism that TrustRank uses to turn a small seed set into a global ranking \cend{17}; AWP's restarts in proportion to sending activity are another instance \cend{3}. Multiplex PageRank lets a node's centrality on one layer shape the walk on another \cend{62}, and the multilayer framework of De Domenico et al.\ lets one walker move within and between the layers of an interconnected network \cend{63}. Both treat layers as different relations over one node set, which matches allowances and transfers among the same addresses. C-PR takes the simplest form of this idea, a per-node weighted sum of two transition matrices with one weight, and S-PR borrows the personalized-teleportation idea, with the allowance score as the teleportation vector (Section~\ref{sub:coupled-graph}).

\dissertationSubheading[sec:account-model]{Accounts, gas, and why logs are the measurement surface}

Ethereum has two kinds of account: externally owned accounts (EOAs) and contract accounts. A private key controls an EOA, whereas a contract account is governed by the code it stores \cend{59}. Every state-changing transaction consumes gas. Logs come from a contract's \texttt{LOG} instructions, which run when its code emits an event (the Solidity \texttt{emit} statement); writing to storage produces no log. A reputation system that means to be auditable must prefer logs to off-chain databases, because logs are a permanent part of the blockchain history that any full node or index provider can recompute independently \cend{60}.

Gas prices also shape the graphs. Cheaper gas on L2 networks densifies $G_{\text{transfer}}$ more than $G_{\text{approve}}$, because transfers recur routinely while an approval is granted once so that later transfers can go through; Chapter~\ref{ch:results} reports the resulting graph sizes.

\dissertationSubheading[sec:approval-security]{Approvals as security-relevant events}

Unlimited approvals (\texttt{type(uint256).max}) save the owner from approving again, at a cost in safety. If the approved account is compromised, every one of the owner's tokens can be stolen until the owner resets the approval. An allowance is therefore a capability and not a record: whoever holds it can move the owner's tokens until it is spent or revoked \cend{18}. EndorseRank interprets a token approval as endorsing one spender's capability to move those tokens. This is why an allowance does not mean the same thing as a transfer, even though both involve the movement of tokens. An account could receive many transfers and never be granted permission to move funds on anyone else's behalf.

In the data, unlimited approvals appear as amounts of $2^{256}-1\approx1.16\times10^{77}$ base units, and they are common (Section~\ref{sub:coupled-graph} gives the counts). EndorseRank's value weight counts them like any other approval, while in the raw allowance layer of C-PR they take almost all of their owner's row.

\dissertationSubheading[sec:awp-in-depth]{AWP in depth: what the transfer-graph baseline measures}

Adaptive Weighted PageRank (AWP), proposed by Do, Do, and Nguyen\cend{3}, ranks the addresses of a transaction graph. Each transaction adds to its edge a logistic function of its time multiplied by a bounded function of its value, so that no transaction adds more than one to the edge weight. The walk restarts at each address in proportion to its activeness, the sum over the addresses it has sent to of the time weight of its latest transaction to each. Ranking accuracy is then validated against inbound degree and inbound value, taken as proxies for the economic attention a wallet attracts. The form of the decay (Equation~\eqref{eq:logistic-decay}) and the in-degree/in-value proxies in Chapters~\ref{ch:methodology} and~\ref{ch:results} are taken from the AWP paper. The paper gives no parameter values; Section~\ref{sub:awp-graph} gives the values used here. The approach has the advantage of being entirely on-chain and rank-based, and it needs no off-chain identity of any kind.

EndorseRank keeps the same template, AWP's edge weights under PageRank plus a rank-correlation check, and changes what the edges mean. In Chapter~\ref{ch:results} we test whether allowance edges yield a score that is construct-distinct, allowance-aligned, and cheaper, and whether scores taken at $t_1$ can predict new approvals by $t_2$.

\dissertationSubheading[sec:method-comparison-table]{Comparison of graph reputation methods}

Table~\ref{tab:method-comparison-lit} compares EndorseRank with the other approaches covered in this chapter by edge semantics, temporal model and validation design.

\begin{table}[htbp]
\centering
\caption{Comparison of selected graph reputation methods (conceptual).}
\label{tab:method-comparison-lit}
\fitwidth{%
\begin{tabular}{p{0.18\linewidth}p{0.24\linewidth}p{0.24\linewidth}p{0.28\linewidth}}
\toprule
Method & Edge semantics & Temporal model & Typical validation \\
\midrule
PageRank (web) \cend{15} & Hyperlinks & None / static snapshot & Retrieval quality \\
EigenTrust \cend{16} & Local trust ratings & Iterative aggregation & P2P file authenticity \\
TrustRank \cend{17} & Links from seed goods & Static / seeded & Spam demotion \\
PageRank\slash HodgeRank on Ethereum \cend{4} & Token transfers & Static window & Not reproduced in this thesis \\
AWP \cend{3} & Transactions, bounded value weights & Logistic decay of edges and restarts & In-degree / in-value \\
RiskProp \cend{12} & Risk-bearing links & Propagation of risk & Account risk labels \\
Multiplex / multilayer PageRank \cend{62,63} & Several relations over one node set & Static layers & Versatility of top-ranked nodes \\
Human Passport \cend{64} & No edges; weighted credentials about the address holder, on-chain and off-chain & Current credential set & Admission threshold; known human and Sybil addresses \\
EndorseRank (this study) & Latest allowances, AWP's weights & Logistic decay of each latest approval & Allowance checks; $t_2$ new approvals \\
C-PR / S-PR (this study) & Allowances mixed with, or seeding, transfers; raw amounts & Allowance snapshot and logistic decay of transfers & Holdout labels; rules fixed in advance \\
\bottomrule
\end{tabular}
}
\end{table}

As Table~\ref{tab:method-comparison-lit} indicates, EndorseRank's contribution lies in its choice of the latest-allowance endorsement edge, and not in the ranking algorithm. On the operator side, our contribution is to couple the two relations into one walk, with a seeded ablation to measure how much this coupling contributes. The coupling is tested out of window, against both constructs, and under a pre-defined rule.

\dissertationSubheading[sec:literature-summary-table]{Summary of the reviewed literature}

Table~\ref{tab:literature-summary} collects in one place the sources underlying this study. Each row records a source's methodology, the strengths this thesis draws on, the limitations it tries to address, and how the source relates to this work. The research gap stated in Section~\ref{sec:research-gap} is grounded in the evidence collected in this table.

{\footnotesize
\setlength{\tabcolsep}{4pt}
\begin{longtable}{p{0.15\linewidth}p{0.20\linewidth}p{0.19\linewidth}p{0.19\linewidth}p{0.18\linewidth}}
\caption{Summary of the reviewed literature: methodology, strengths, limitations, and relation to this study.}
\label{tab:literature-summary}\\
\toprule
Source (year) & Methodology & Strengths & Limitations & Relation to this study \\
\midrule
\endfirsthead
\multicolumn{5}{l}{\small Table~\thetable\ (continued)}\\
\toprule
Source (year) & Methodology & Strengths & Limitations & Relation to this study \\
\midrule
\endhead
\midrule
\multicolumn{5}{r}{\small continued on next page}\\
\endfoot
\bottomrule
\endlastfoot
Page et al.\ (1998) \cend{15}; Langville and Meyer (2006) \cend{23} & Eigenvector ranking of a hyperlink graph, solved by power iteration & Interpretable propagation; $O(|E|)$ per iteration; damping yields a unique Perron vector & Edges can only mean hyperlinks; no notion of time & Same solver and complexity argument, used unchanged for both methods compared \\
Kamvar et al.\ (2003) \cend{16}; Gy\"{o}ngyi et al.\ (2004) \cend{17} & Aggregation of trust in P2P networks; propagation from seeds to demote spam & Normalization and seed sets curb inflation; the adversary is modeled explicitly & Depends on explicit ratings or curated seeds; not applied on-chain & Informs the Sybil-cost discussion and the uniform-teleportation baseline \\
Do and Do (2023) \cend{4} & PageRank and HodgeRank applied to Ethereum transaction graphs as a measure of social credit & Treats transfer connectivity as relevant to credit; data are fully public & Transfers conflate payment with trust; allowance semantics absent & Earlier transfer-graph work, kept separate from AWP in this thesis \\
Do, Do, and Nguyen (2023), AWP \cend{3} & Transaction PageRank with bounded value weights, logistic recency and restarts weighted by sending activity; rank correlation with in-degree/in-value as validation & Validates on ranks, so no default labels are needed; recency modeled explicitly & Proxies are computed from the same transfer data that forms the graph (mechanical alignment); dense graph; no parameter values reported & Baseline, used as published (Section~\ref{sub:awp-graph}); the decay parameters and $b$ are set here, and EndorseRank takes the same edge weights \\
Lin et al.\ (2025), RiskProp \cend{12} & Network propagation of account risk, paired with a de-anonymization score & Structure shown to add information beyond activity counts; evaluated with labels & Needs risk labels; aimed at fraud, not trust or authorization & Argues for finer edge semantics; a simplified liquidation-weighted transfer walk is archived (Appendix~\ref{ch:appendix-archive}) \\
Wu et al.\ (2022) \cend{24} & Phishing detection through network embedding of Ethereum transaction graphs & Graph features outperform per-account features; large labeled dataset & Supervised; its labels mark phishing, not creditworthiness & Backs the premise that graph position is informative; not a competing method \\
Hassija et al.\ (2020) \cend{13}; Kotb (2023) \cend{14} & Blockchain-data credit scoring using prospect theory and machine learning & Aimed directly at lending decisions & Need labels or off-chain features; only partly interpretable & A contrasting class; EndorseRank remains within auditable graph propagation \\
Bellini et al.\ (2020) \cend{25}; Almasoud et al.\ (2020) \cend{26} & Surveys of blockchain-based trust and reputation systems & Chart the design space and find explicit ratings and volume-based signals dominant & Authorization edges left unevaluated; nothing measured on L2 & Show that none of the surveyed designs uses latest-allowance edges \\
Victor (2020) \cend{40}; Victor and Weintraud (2021) \cend{28}; Cong et al.\ (2023) \cend{29} & Heuristics for clustering addresses; detection of wash trading on DEXs and exchanges & Measure the share of on-chain activity that is artificial or multi-address & No reputation score proposed & Show why transfer-based signals are open to inflation, which argues for the allowance construct \\
Cronbach and Meehl (1955) \cend{22}; Messick (1995) \cend{33} & Psychometric construct validity & Distinguish what a score is designed to measure from external criteria & Not specific to blockchains & Basis for treating the contemporaneous checks as construct checks and the temporal holdouts as out-of-window evidence \\
Human Passport (n.d.) \cend{64}; Siddarth et al.\ (2020) \cend{65}; Bordeianu and Popescu (2026) \cend{66} & Weighted credentials about an address holder, on-chain and off-chain, against a passing threshold; reviews of proof-of-personhood defenses & Deployed at scale; uses off-chain evidence against Sybil attacks & A gate and not a graded ordering; relations between addresses play no part; credential farming remains & Shows how trust is established in practice, and what EndorseRank does not offer: no personhood guarantee \\
Efron (1979) \cend{34}; DiCiccio and Efron (1996) \cend{35} & Bootstrap resampling and bootstrap-based confidence intervals & Uncertainty for rank statistics with no distributional assumption; paired contrasts & Units must be exchangeable & Source of the 95\% intervals and $\Delta\tau$ contrasts reported in Chapter~\ref{ch:results} \\
\end{longtable}
}

\dissertationSubheading[sec:research-gap]{Research gap}

Five gaps emerge from Table~\ref{tab:literature-summary} and Section~\ref{sec:related-measurement}. First, every on-chain reputation method discussed here (PageRank/HodgeRank \cend{4}, AWP \cend{3}, RiskProp \cend{12}, and the designs in the surveys \cend{25}) builds its graph from transfers or explicit ratings, and the attestation tooling deployed in DeFi scores the holder of an address without any graph of relations between addresses \cend{64}. None uses the ERC-20 allowance, the on-chain record of a deliberate delegation of spending power that every ERC-20 token carries, and it remains unknown whether a latest-allowance endorsement graph yields a distinct and meaningful ordering (O1 and RQ1). A second gap lies in validation. The transfer-graph baseline checks its scores against in-degree and in-value computed from the same transfer data used to build the graph, so some agreement is unavoidable; no report applies this procedure to an allowance construct and its transfer counterpart with confidence intervals and paired difference tests on the same set of addresses (O2 and RQ2). Third, runtimes for blockchain PageRank are often published without edge counts, memory use or iteration counts under a fixed protocol, which makes it hard to attribute any efficiency claim to sparsity as opposed to implementation (O3 and RQ3). The fourth gap is methodological and has to do with time. Graph statistics from the same window are often presented as validation, although such scores cannot be shown to predict future behavior; the out-of-window check used here is instead a temporal holdout with a future-tense label for each construct and raw degree at the time of the freeze as the baseline (O4 and RQ4). Finally, multilayer ranking has been developed for social, transport and citation networks \cend{62,63} but has not been used on two on-chain relations over one set of addresses. No study reports whether combining an authorization layer with a payment layer improves on the payment layer alone, or beats both layers, under rules fixed before the outcome is seen (O5 and RQ5).

\dissertationSubheading[sec:open-problems]{Open problems motivating this research}

The literature leaves several problems in on-chain reputation unresolved, and these motivate the present study:

\begin{enumerate}
\item Edge semantics: When identities are pseudonymous, which on-chain relations capture trust or authorization best \cend{2}?
\item Validation without defaults: On what basis can methods be compared when loan-default labels are scarce \cend{4}?
\item Computational cost: As graphs grow, can reputation scores still be refreshed at interactive latencies \cend{23}?
\item Time-split checks: What keeps an evaluation from mistaking same-window graph statistics for out-of-window evidence \cend{22}?
\item Adversarial robustness: How much does it cost to forge reputation when identities are cheap \cend{11}?
\end{enumerate}

EndorseRank takes on (1)--(4) through allowance edges, contemporaneous checks, runtime benchmarks, and a temporal holdout. Problem (5) is analyzed at the model level in Chapter~\ref{ch:discussion}; a full empirical adversarial evaluation is left for future work.

Chapter~\ref{ch:methodology} turns this gap statement into an operational design: the data window and cohorts, the construction of both graphs and their solution with one shared PageRank implementation, the proxy families, the temporal holdout, and the timing and bootstrap protocols behind the results of Chapter~\ref{ch:results}.

